\section{Scénario de Stress et P\&L}
\label{sec:stress}

\subsection{Scénario}

On applique simultanément :
\begin{itemize}
  \item $F_T \to F_T \times 1.10$ (choc spot $+10\%$ en relatif)
  \item $\sigma_{SSVI}(k,t) \to \sigma_{SSVI}(k,t) - 0.10$ (choc vol $-10\,\text{pp}$ absolu)
  \item $T \to T - \tfrac{1}{52}$ (avancement de 1 semaine)
\end{itemize}

\textbf{Convention du choc de vol.}
Un choc absolu de $-10\,\text{pp}$ signifie : si $\sigma = 80\%$, alors
$\sigma_{new} = 70\%$. C'est plus conservateur qu'un choc relatif $\times 0.9$
(qui donnerait $72\%$) et cohérent avec la pratique des risk managers.

\subsection{Calcul du P\&L}

\begin{equation}
  \PnL_{total} = \PnL_{produit} + \PnL_{hedge}
\end{equation}
avec $\PnL_{produit} = -(V_{produit}^{new} - V_{produit}^{init})$
(position vendeuse) et $\PnL_{hedge} = \sum_i q_i (V_i^{new} - V_i^{init})$.

\subsection{Résultats}

% TODO : tableau de revalorisation (F_avant/après, σ_avant/après, prix_avant/après)
% TODO : waterfall P&L (plot_pnl_waterfall)
